\section{Iterator}
\label{sec:pat-iterator}

\problemstatement{Provide a way to access the elements of an aggregate
object sequentially without exposing its underlying representation, so
clients traverse a collection the same way regardless of how it is stored.}

\begin{patternbox}[When You Reach for It]
The trigger is \emph{``traverse a collection without revealing whether
it's an array, list, or tree.''} When you want uniform iteration and the
freedom to change the container's internals, expose an iterator instead of
raw structure.
\end{patternbox}

\subsection*{Understanding the pattern}

Define an iterator interface with \texttt{hasNext()} and \texttt{next()}.
The collection produces an iterator that knows how to walk \emph{its}
particular structure. Client loops depend only on the iterator interface,
so swapping a vector for a linked list -- or offering multiple traversal
orders -- does not touch client code. This is the pattern the C++ STL and
range-based \texttt{for} are built on.

\begin{ideabox}[Externalise Traversal State]
Move the ``where am I in the collection'' cursor out of the client and
into a dedicated iterator object. Multiple independent iterators can then
traverse the same collection at once, and the collection's storage stays
hidden.
\end{ideabox}

\subsection*{C++ implementation}

\begin{lstlisting}
#include <bits/stdc++.h>
using namespace std;

template <typename T>
struct Iterator {
    virtual bool hasNext() const = 0;
    virtual T next() = 0;
    virtual ~Iterator() = default;
};

class IntCollection {                           // hides its storage
    vector<int> data;
public:
    void add(int x) { data.push_back(x); }

    class Iter : public Iterator<int> {
        const vector<int>& ref; size_t pos = 0;
    public:
        Iter(const vector<int>& r) : ref(r) {}
        bool hasNext() const override { return pos < ref.size(); }
        int next() override { return ref[pos++]; }
    };

    unique_ptr<Iterator<int>> iterator() const {
        return make_unique<Iter>(data);
    }
};

int main() {
    IntCollection c;
    c.add(10); c.add(20); c.add(30);
    auto it = c.iterator();
    while (it->hasNext()) cout << it->next() << " ";   // storage hidden
    cout << "\n";
    return 0;
}
\end{lstlisting}

\begin{complexitybox}[Trade-offs]
\textbf{Pros.} Uniform traversal; hides representation; supports multiple
simultaneous and multiple-order traversals; simplifies the collection's
own interface. \textbf{Cons.} Extra objects for simple cases; an iterator
can be invalidated if the collection is modified during traversal.
\end{complexitybox}

\begin{takeawaybox}
Iterator externalises traversal into an object with
\texttt{hasNext}/\texttt{next}, decoupling clients from a collection's
internal structure. It is so fundamental that most languages bake it in
(STL iterators, \texttt{for-each}). Mention iterator invalidation as the
classic pitfall.
\end{takeawaybox}
